\begin{savequote}
\includegraphics[width=0.7\linewidth]{./images/noauth/AASdev20090228-img110.jpg}
\end{savequote}

\chapter
[\CO{[KHD]} Einführung in den Gro{\ss}schadens- und Katastrophenhilfsdienst]
{\CC{[KHD]} Einführung in den Gro{\ss}schadens- und Katastrophenhilfsdienst}
\label{chp:KHD}
% \AasRsN{RS~XI}

\begin{inparadesc}
\item[Maintainer:] Roman Koch
\item[Reviewer:] Gerald Höritzmiller
\end{inparadesc}
% \emph{Maintainer: Roman Koch}

% Changelog:
% 
% 2009-02-20: Fertig für Kurs 2009/02.
% 
% 2009-02-25: Start Changelog.
% 
%\write18{git log > version.log}

%\input{|"git log"}

% \A{\tiny
% \input{include-khd-chapter.commit20.log} 
% }

 
\minitoc 

\section{Katastrophen, Großschadensereignisse, Unfall}
% TODO : Layout-Anpassung
\TODO{GABS}{yyyy-mm-dd}{Layout-Anpassung notwendig.}{Dieser Teil entspricht nicht dem Standard-Layout!}


\begin{minipage}{\linewidth}
\begin{center}

\includegraphics[width=\linewidth]{./images/noauth/AASdev20090228-img109.jpg}
 
\end{center}

\Parallel{
\includegraphics[width=\linewidth]{./images/noauth/AASdev20090228-img110.jpg}
}{
\Captionof{figure}{ICE-Unglück von Eschede am 3. Juni
1998. \SourcePicture{Wikipedia, Nils Fretwurst; Lizenz: gemein frei
({\quotedblbase}public domain{\textquotedblleft})}}

Verstorbene:~101

Schwerverletzte:~88  

Leicht- und Unverletzte:~106

}

\end{minipage}






\subsection{Grundlegende Begriffe}






\subsubsection{Unfall (Bagatelle, Notfall)}

\Layout{0}{\TxBeschreibung}{%
Der Unfall ist plötzlich und unerwartet, die Dauer
des Einsatzes wird in Minuten bis 1 Std. gemessen. Die
\E{Infrastruktur}\footnote{\emph{Infrastruktur:} Wichtige Basiseinrichtungen und
Dienstleistungen, die das öffentliche Leben ermöglichen; z.\,B. Straßen,
Elektrizität, Telefon, Krankenhäuser, Verwaltungseinrichtungen, Nahrungs- und
Wasserversorgung, Abfall- und Abwasserentsorgung, usw.} ist \E{intakt}.
Die Anzahl der Patienten übersteigt nicht  14 Personen (Richtwert).

Die Erstversorgung erfolgt innerhalb von Minuten durch Rettungsdienst nach
den Regeln der \E{Individualmedizin}. Der Abtransport wird durch Routinebetrieb
mit örtlich vorhandenen Mitteln durchgeführt. Die Definitive Versorgung erfolgt 
in umliegenden Krankenhäusern, ebenfalls nach den Regeln der
\E{Individualmedizin}. Das aufgewendete Material und das Personal stammt aus dem
Normalbetrieb des Rettungs- und Krankentransportes.
}{\prettyref{tab:uebersicht-unfall-grossunfall-katastrophe}}{}{}{}



% 
% \begin{tabulary}{\linewidth}{LL}
% \toprule
% \TabSub{Ereignis:}
%  &
% plötzlich, unerwartet
% \\\addlinespace
% \TabSub{Dauer:}
%  &
% Minuten bis 1 Std.
% \\\addlinespace
% \TabSub{Infrastruktur:}\MarginBugfix{\emph{Infrastruktur:} Wichtige
% Basiseinrichtungen und Dienstleistungen, die das öffentliche Leben ermöglichen;
% z.\,B. Straßen, Elektrizität, Telefon, Krankenhäuser, Verwaltungseinrichtungen, 
% Nahrungs- und Wasserversorgung, Abfall- und Abwasserentsorgung, usw.} 
% &
% intakt 
% \\\addlinespace
% \TabSub{Verletzte:}
%  &
% bis 14
% \\\addlinespace
% \TabSub{Erstversorgung:}
%  &
% innerhalb von Minuten durch Rettungsdienst nach Regeln der
% Individual\-medizin
% \\\addlinespace
% \TabSub{Abtransport:}
%  &
% unbehindert Routinebetrieb mit örtlich vorhandenen Mitteln
% \\\addlinespace
% \TabSub{Definitive Versorgung:}
%  &
% in umliegende Krankenhäuser nach Regeln der Individualmedizin
% \\\addlinespace
% \TabSub{Material und Mannschaft:}
%  &
% Normalbetrieb Rettungs- und Krankentransport
% \\\bottomrule
% \end{tabulary}


\subsubsection{Großunfall (Großschadensereignis)}
\index{Großunfall}
\index{Großschadensereignis}
\label{topic:grossschaden}

\Definition{{\quotedblbase}Ein Großunfall liegt vor, wenn anzunehmen ist,
dass das Ereignis mit den örtlichen personellen und materiellen
Kräften und Mitteln nicht bewältigt werden kann, aber keine
erklärte Katastrophensituation vorliegt{\textquotedblleft}}

\Layout{0}{\TxBeschreibung}{
Das Ereignis ist \E{lokal begrenzt} und dauert mehrere Stunden, die
\E{Infrastruktur} ist jedoch \E{intakt}. Die Erstversorgung von 15
oder mehr Patienten (Richtwert) erfolgt innerhalb von Minuten bis Stunden durch
den Rettungsdienst nach den Regeln der \E{Katastrophenmedizin}. Der Abtransport
erfolgt relativ unbehindert nach der Erstversorgung, es sind genügend Transportmittel vorhanden.

Die Definitive Versorgung erfolgt in umliegende Krankenhäuser nach den Regeln
der \E{Individualmedizin}. Das Personal setzt sich aus alle verfügbaren Helfern
zusammen, welche sowohl aus dem Normalbetrieb des Rettungs- und Krankentransportes und von Sondereinsatzgruppen stammt. Es
wird zusätzliches Material eingesetzt.
}{\prettyref{tab:uebersicht-unfall-grossunfall-katastrophe}}{}{}{}


% \begin{tabulary}{\linewidth}{LL}
% \toprule
% \TabSub{Ereignis:}
%  &
% lokal begrenzt
% \\\addlinespace
% \TabSub{Dauer:}
%  &
% mehrere Stunden
% \\\addlinespace
% \TabSub{Infrastruktur:}
%  &
% intakt
% \\\addlinespace
% \TabSub{Verletzte:}
%  &
% ab 15
% \\\addlinespace
% \TabSub{Erstversorgung:}
%  &
% innerhalb von Minuten bis Stunden durch den Rettungsdienst nach den
% Regeln der Katastrophenmedizin
% \\\addlinespace
% \TabSub{Abtransport:}
%  &
% relativ unbehindert nach der Erstversorgung, genügend vorhandene
% Transportmittel
% \\\addlinespace
% \TabSub{Definitive Versorgung:}
%  &
% in umliegende Krankenhäuser nach den Regeln der Individualmedizin
% \\\addlinespace
% \TabSub{Material und Mannschaft:}
%  &
% alle verfügbaren Helfer, Normalbetrieb Rettungs- und Krankentransport
% inkl. Sondereinsatzgruppe und zusätzlichem Material
% \\\bottomrule
% \end{tabulary}


\subsubsection{Katastrophe (außergewöhnliche Lage)}
\index{Katastrophe}\label{topic:katastrophe}
% \noindent \Gls{katastrophe} \Lang{griech.}: schweres Unglück,
% Zusammenbruch, Wendung zum Schlimmen \cite{BertelsmannVolkslexikon25,BrockhausDrei2Bd2}.
\Lang{griech.}: schweres Unglück, Zusammenbruch,
  Wendung zum Schlimmen \cite{BertelsmannVolkslexikon25,BrockhausDrei2Bd2}.

\Definition{Ein \Es{unvorhergesehenes Ereignis} mit einem
\Es{außergewöhnlichen Schadensausmaß} das unmittelbar bevorsteht oder
bereits eingetreten ist, eine \Es{konkrete Gefahr für Menschen}\Z{, Tiere,
Umwelt, Kulturgüter} und Sachwerte sowie für die \Es{Infrastruktur} zur
Sicherstellung der Versorgung mit lebensnotwendigen Gütern und Dienstleistungen, welches
der \Es{koordinierten Führung durch die Behörde} bedarf. }

\Layout{0}{\TxBeschreibung}{
Das Ereignis betrifft ein \E{(über)-regionales Gebiet} und kann
\E{Tage bis Wochen} andauern. Die \E{Infrastruktur} ist \E{gestört}, oder sogar
zerstört. Die Anzahl der Verletzten variiert je nach Art
des Ereignisses. Die Erstversorgung erfolgt innerhalb von Stunden bis Tagen
durch den KHD\footnote{KHD: Katastrophen- und
Großschadenshilfsdienst}-Einheiten nach den Regeln der \E{Katastrophenmedizin}.
Der Abtransport ist stark behindert und verzögert, es gibt nur wenige oder keine
vorhandene Transportmittel. 

Die Definitive Versorgung erfolgt vorwiegend auf Verbandsplätzen, in
Behelfsspitälern, oder in Lazaretten nach den Regeln der
\E{Katastrophenmedizin}. Das Material und und Personal setzt sich zusammen aus
Überlebenden und verfügbaren Helfer, einsatzbereiten Fahrzeugen,
Kat-Einsatz aller Hilfsorganisationen, überregionaler und
internationaler Hilfe.
}{\prettyref{tab:uebersicht-unfall-grossunfall-katastrophe}}{}{}{}

\Layout{0}{Arten von Katastrophen}{
Es gibt verschiedene Katastrophenarten
\begin{center}
\begin{tabularx}{\linewidth}{XX}
\T{Naturkatastrophen}
 &
Erdbeben,\"Uberschwemmungen, Erdrutsche/Lawinen etc.
\\
\T{Technische Katastrophen}
 &
Flugzeugabsturz, Eisenbahnunglück, Reaktor\-unfall, Schiffs\-untergang
etc.
\\
\T{Sekundärkatastrophen}
 &
Hungersnot, Seuchen etc.
\\
\end{tabularx}
\end{center}
}{
\begin{itemize}
  \item Naturkatastrophen.
  \item Technische Katastrophen.
  \item Sekundärkatastrophen.
\end{itemize}
}{}{}{}


% \begin{center}
% \begin{tabular}{m{3.31cm}m{9.99cm}}
% \toprule
% \TabSub{Ereignis:}
%  &
% (über)-regionales Gebiet
% \\\addlinespace
% \TabSub{Dauer:}
%  &
% Tage, Wochen
% \\\addlinespace
% \TabSub{Infrastruktur:}
%  &
% gestört, zerstört
% \\\addlinespace
% \TabSub{Verletzte:}
%  &
% keine bis viele, je nach Art des Ereignisses
% \\\addlinespace
% \TabSub{Erstversorgung:}
%  &
% innerhalb von Stunden bis Tagen durch KHD\footnotemark{}{}-Einheiten
% nach den Regeln der Katastrophenmedizin
% \\\addlinespace
% \TabSub{Abtransport:}
%  &
% stark behindert verzögert, wenige/nicht vorhandene Transportmittel
% \\\addlinespace
% \TabSub{Definitive Versorgung:}
%  &
% auf Verbandsplätzen in Behelfsspitälern, in Lazaretten nach Regeln
% der Katastrophenmedizin
% \\\addlinespace
% \TabSub{Material und Mannschaft:}
%  &
% \"Uberlebende und verfügbare Helfer, einsatzbereite Fahrzeug,
% Kat-Einsatz aller Hilfsorganisationen, überregionale und
% internationale Hilfe
% \\\bottomrule
% \end{tabular}
% \end{center}
% \footnotetext{KHD: Katastrophenhilfsdienst}

\begin{table}[H]
\begin{gWideParEnv}
\Captionof{table}{Übersicht
Unfall\,--\,Grossunfall\,--\,Katastrophe\label{tab:uebersicht-unfall-grossunfall-katastrophe}}
\begin{tabularx}{\linewidth}{p{8em}XXX}
&
\TabHeading{Unfall}
&
\TabHeading{Großunfall}
&
\TabHeading{Katastrophe}
\\\addlinespace\midrule\addlinespace
\TabSub{Ereignis:}
 &
plötzlich, unerwartet
&
lokal begrenzt
&
(über)-regionales Gebiet
\\\addlinespace
\TabSub{Dauer:}
 &
Minuten bis 1 Std.
&
mehrere Stunden
&
Tage, Wochen
\\\addlinespace
\TabSub{Infrastruktur:}

&
intakt 
&
intakt
&
gestört, zerstört
\\\addlinespace
\TabSub{Patienten:}
 &
{\textless} 15 (Richtwert)
&
$\geq$ 15 (Richtwert)
&
keine bis viele, je nach Art des Ereignisses
\\\addlinespace
\TabSub{Erstversorgung:}
 &
innerhalb von Minuten durch Rettungsdienst nach Regeln der
Individual\-medizin
&
innerhalb von Minuten bis Stunden durch den Rettungsdienst nach den
Regeln der Katastrophenmedizin
&
innerhalb von Stunden bis Tagen durch KHD\footnotemark{}{}-Einheiten
nach den Regeln der Katastrophenmedizin
\\\addlinespace
\TabSub{Abtransport:}
 &
unbehindert Routinebetrieb mit örtlich vorhandenen Mitteln
&
relativ unbehindert nach der Erstversorgung, genügend vorhandene
Transportmittel
&
stark behindert verzögert, wenige/nicht vorhandene Transportmittel
\\\addlinespace
\TabSub{Definitive Versorgung:}
 &
in umliegende Krankenhäuser nach Regeln der Individualmedizin
&
in umliegende Krankenhäuser nach den Regeln der Individualmedizin
&
auf Verbandsplätzen in Behelfsspitälern, in Lazaretten nach Regeln
der Katastrophenmedizin
\\\addlinespace
\TabSub{Material und Mannschaft:}
 &
Normalbetrieb Rettungs- und Krankentransport
&
alle verfügbaren Helfer, Normalbetrieb Rettungs- und Krankentransport
inkl. Sondereinsatzgruppe und zusätzlichem Material
&
\"Uberlebende und verfügbare Helfer, einsatzbereite Fahrzeug,
Kat-Einsatz aller Hilfsorganisationen, überregionale und
internationale Hilfe
\\\bottomrule
\end{tabularx}
\footnotetext{KHD: Katastrophenhilfsdienst}
\end{gWideParEnv}
\end{table}



\subsubsection{Individualmedizin}
\index{Individualmedizin}

\Definition{Bedarfsgerechte, maximale und bestmögliche Versorgung jedes
Patienten.}

\Layout{0}{\TxBeschreibung}{
Es besteht freie Arztwahl, welche auf das Vertrauensverhältnis zwischen
Patient und Arzt ausgerichtet ist.
}{
\begin{itemize}
  \item Bedarfsgerechte, maximal- und bestmögliche Versorgung jedes
Patienten, wobei örtlich vorhandene Ressourcen für den Einzelnen ausgeschöpft
werden.
\end{itemize}
}{}{}{}



\subsubsection{Katastrophenmedizin}
\index{Katastrophenmedizin}

\Definition{Massenversorgung von Verletzten und
Erkrankten mit beschränkten Mitteln und dem Zwang zur Einteilung nach
\underline{Dringlichkeit} und \underline{Möglichkeit} zur Behandlung.}

\Layout{0}{\TxBeschreibung}{
Eine Vielzahl von Hilfsbedürftigen stehen wenigen Ärzten,
Sanitätern und sonstigen Helfern mit oft nicht ausreichenden Mitteln
gegenüber. Ziel muss es sein, unter den gegebenen Umständen so
viele Leben wie möglich zu retten, \Es{mit bewusstem Verzicht auf zeit-
und personalaufwendige Maximalversorgung des Einzelnen zugunsten der
lebensrettenden Minimalversorgung Vieler.}

Ein wesentliches Element der Katastrophenmedizin ist die Einteilung der
Patienten gemäß der Dringlichkeit und Druchführbarkeit ihrer Behandlung oder
ihres Transportes. Diese Einteilung nennt man Triage.
}{
\begin{itemize}
  \item Aus beschränkten Mitteln das Beste für Viele herausholen, wobei auf
  maximale Versorgung Einzelner bewusst verzichtet werden muss.
\end{itemize}
}{}{}{}




















\subsubsection{Triage und Triagegruppen}
\label{topic:triage}
\index{Triage}

\Definition{\Es{Begutachten} aller Patienten mit geringen Mitteln und
  Aufwand sowie das \Es{Einteilen} in Triagegruppen gemäß der
  \Es{Behandlungsdringlichkeit} und \Es{-möglichkeit} mit dem Ziel mit
  den vorhandenen Resourcen soviel als möglich zu erreichen. \Z{Triage
  \Lang{franz.}: Sichtung, Auswahl}}

\Fazit{Grundsatz: 

{\quotedblbase}Life before Limbs{\textquotedblleft} = Leben vor
Gliedmaßen, Funktion vor Schönheit!
}

\noindent\prettyref{tab:khd-behandlungsstellen-uebersicht} gibt eine Übersicht
über die Triagegruppen.

\subsubsection{Die Transportpriorität ist unabhängig von der Triagegruppe}
\label{topic:transportprioritaet}
\index{Transportpriorität}

Nach einer erfolgten Behandlung kann eine Transportpriorität zugewiesen
werden. Man unterscheidet hierbei: 

\begin{description}
\item[A:] Hohe Transportpriorität: Der rasche Transport zur
frühzeitigen Fachbehandlung nach ärztlicher Notversorgung ist
angesagt. 
\item[B:] Niedrige Transportpriorität: Der Transport zur Fachbehandlung
kann verzögert erfolgen. 
\end{description}


\begin{table}[H]
\small
\noindent\begin{gWideParEnv}
\Captionof{table}{Triagegruppen,
Übersicht.\label{tab:khd-behandlungsstellen-uebersicht}\label{tab:khd-triagegruppen-uebersicht}}
\noindent\begin{tabularx}{\linewidth}{XXXX}
\TabHeading{I -- Sofortbehandlung vor Ort}
&
\TabHeading{II -- Dringliche Behandlung}
&
\TabHeading{III -- Warten Leichtverletzte}
&
\TabHeading{IV -- Warten Schwerverletzte}
\\\addlinespace\midrule\addlinespace
Lebensrettende Soforteingriffe
&
Schwerverletzte
&
Leichtverletzte
&
Schwerverletzte
\\\addlinespace
\includegraphics[width=3cm]{./images/noauth/sanhist-triage-1.png}
&
\includegraphics[width=3cm]{./images/noauth/sanhist-triage-2.png}
&
\includegraphics[width=3cm]{./images/noauth/sanhist-triage-3.png}
&
\includegraphics[width=3cm]{./images/noauth/sanhist-triage-4.png} 
\\\addlinespace\midrule\addlinespace
\multicolumn{4}{c}{\TabSub{Welche Patienten?}}
\\\addlinespace\midrule\addlinespace
Patienten, welche ohne sofortige Behandlung sterben bzw. schwere
Schäden erleiden
&
Behandlung und Herstellung der Transportfähigkeit von schwer
verletzten/erkrankten Personen, eine weitere Behandlung ist nur mehr im
Krankenhaus möglich.
&
Patienten, welche nur einer Minimalbehandlung bedürfen.
&
Patienten, welche mit dem zur Verfügung stehenden Mitteln
(Personal, Material) \Es{derzeit nicht} versorgt werden können. Abwartende
Behandlung.
\\\addlinespace\midrule\addlinespace
\multicolumn{4}{c}{\TabSub{Zum Beispiel}}
\\\addlinespace\midrule\addlinespace
\begin{itemize}
    \item[--] Atemstörungen
    \item[--] Spannungs\-pneu\-mo\-thorax
    \item[--] schwere äußere Blutungen
    \item[--] schwerer traumatischer Schock
\end{itemize}
&
\begin{itemize}
    \item[--] SHT
    \item[--] Wirbelverletzungen mit Lähmungen
    \item[--] Bauchtrauma, innere Blutungen
    \item[--] schwere Thoraxtraumen
    \item[--] große Gefäßverletzungen
    \item[--] offene Knochenbrüche und  Gelenks\-verletzungen
    \item[--] Augenverletzungen
    \item[--] große Weichteilverletzungen
    \item[--] Verbrennungen 15 - 30\%
    \item[--] geschlossene Knochenbrüche und Gelenks\-ver\-letzungen
\end{itemize}
&
\begin{itemize}
    \item[--] Verbrennungen bis 15\%
    \item[--] kleine Weichteilwunden
    \item[--] einfache Knochenbrüche
    \item[--] Prellungen, Zerrungen
    \item[--] leichtes SHT
\end{itemize}
&
\begin{itemize}
    \item[--] schweres Polytrauma
    \item[--] schwerstes \Abk{SHT}
    \item[--] offene Körperhöhlentraumen
    \item[--] Mehrhöhlenverletzungen
    \item[--] Verbrennungen über 50\%
\end{itemize}
\\\addlinespace\midrule\addlinespace
\multicolumn{4}{c}{\TabSub{Maßnahmen, Beispiele}}
\\\addlinespace\midrule\addlinespace
\begin{itemize}
    \item[--] \"Uberprüfung der Triage
    \item[--] Volumenersatz
    \item[--] Intubation
    \item[--] Thoraxdrainage
    \item[--] Noteingriffe
\end{itemize}
&
\begin{itemize}
    \item[--] Verbände
    \item[--] Schockbehandlung
    \item[--] Fixierung/Schienung
    \item[--] Schmerzbekämpfung
    \item[--] pflegerische Maßnahmen
    \item[--] Lebensrettende Sofortmaßnahmen
\end{itemize}
&
\begin{itemize}
    \item[--] \"Uberprüfung der Triage
    \item[--] pflegerische Maßnahmen
    \item[--] Betreuung (Kriseninterventionsteam (\Abk{KIT})
    \item[--] Entlassung ambulant Behandelter zur "`Sammelstelle Unverletzte"'
\end{itemize}
&
\begin{itemize}
    \item[--] \"Uberprüfung der Triage
    \item[--] Behandlung  (Schmerzbekämpfung)
    \item[--] Psychologische Betreuung, evtl. durch \Abk{KIT}
    \item[--] \"Uberwachung
\end{itemize}
\\\addlinespace\bottomrule

\end{tabularx}
\end{gWideParEnv}
\end{table}




\clearpage
\subsection{Regeln und Ziele bei Katastrophen und Großschadenslagen}

\subsubsection{Bei Großunfällen und Katastrophen versucht man das Beste
herauszuholen}
\label{topic:katastrophe-material}

\Layout{0}{\TxBeschreibung}{
Ziel ist es, das \Es{Bestmögliche} für die größte Anzahl zur rechten Zeit am
rechten Ort zu erreichen. Das ist nur erreichbar wenn  der \E{Bestgeeignete} auf
jeder Stufe führt, klare und eindeutige \Es{Führungsstrukturen} vorhanden sind,
alle Einsatzkräfte ihre \E{Aufgabe kennen},  die \Es{Führungspyramide} von allen
strikt eingehalten wird, das dazu benötigte \E{Material vorhanden},
bereitgestellt und rasch einsatzbereit ist, die Organisation eingespielt
und eingeübt ist sowie Einsätze -- soweit möglich -- geplant und vorbereitet
sind . 

Dafür wird zusätzliche Sonderausrüstung erforderlich, da ein
erhöhter Bedarf an Material, Medikamenten und Personal besteht.  Wir müssen mit mit einem
großen Patientenanfall,  besonderen Verletzungsmustern, einer längeren
Einsatzzeit vor Ort und nicht zuletzt mit einer Überforderung der Helfer rechnen.
}{
\TabSub{Ziele}
\begin{itemize}
    \item Das Bestmögliche, für die größte Anzahl,  zur rechten Zeit, am
    rechten Ort.
\end{itemize}

\TabSub{Voraussetzungen}
\begin{itemize}
    \item Führung durch Bestgeeigneten
    \item Führungsstrukturen, -pyramide
    \item Kenntnis der Aufgaben
    \item Material vorhanden und EB
    \item Übung, Plan und Vorbereitung
\end{itemize}

\TabSub{Sonderausrüstung} da
erhöhter Bedarf an:
\begin{itemize}
    \item Material
    \item Medikamenten
    \item Personal
\end{itemize}

\TabSub{Wir müssen rechnen mit}
\begin{itemize}
    \item großem Patientenanfall
    \item besondere Verletzungsmuster
    \item längere Einsatzzeit
    \item Überforderung der Helfer
\end{itemize}
}{}{}{}



% 
% \begin{tabularx}{\linewidth}{XX}
% \hline
% \TabHeading{Ziele}
% 
% \begin{itemize}
%     \item das Bestmögliche
%     \item für die grö{\ss}te Anzahl
%     \item zur rechten Zeit
%     \item am rechten Ort
% \end{itemize}
%  &
% \TabHeading{Das ist nur erreichbar wenn:}
% 
% \begin{itemize}
%     \item der Bestgeeignete auf jeder Stufe führt
%     \item alle Einsatzkräfte Ihre Aufgabe kennen
%     \item das dazu benötigte Material vorhanden, bereitgestellt und rasch
%     einsatzbereit ist
%     \item klare und eindeutige Führungsstrukturen vorhanden sind
%     \item die Führungspyramide von allen strikt eingehalten wird
%     \item die Organisation eingespielt/eingeübt ist
%     \item Einsätze geplant und vorbereitet sind
% \end{itemize}
% \\\hline
% \TabHeading{Dafür werden Sonderausrüstungen zusätzlich erforderlich,} da
% ein erhöhter Bedarf besteht an:
% 
% \begin{itemize}
%     \item Material
%     \item Medikamenten
%     \item Personal
% \end{itemize}
% &
% \TabHeading{Wir müssen rechnen mit}
% \begin{itemize}
%     \item einem gro{\ss}en Patientenanfall
%     \item besonderen Verletzungsmustern
%     \item einer längeren Einsatzzeit vor Ort
%     \item Überforderung der Helfer
% \end{itemize}
% \\\hline
% \end{tabularx}



\subsubsection{Katastrophenhilfegesetze regeln die Zuständigkeiten}
\label{topic:katastrophenhilfsgesetze}

Die Katastrophenhilfsgesetze regeln die Aufgaben der Länder,
Organisationen und Zuständigkeiten in Ihrem Bereich.
\Es{Ausschließlich eine Behörde erklärt ein außergewöhnliches
Schadensereignis zur Katastrophe} im Sinne des Gesetzes. Behördliche
Alarmpläne dienen zur Festlegung der Maßnahmen der
Katastrophenhilfe und zur Koordination der im Anlassfall tätig
werdenden Einrichtungen. Je nach Ausdehnung des Schadensgebietes gibt
es unterschiedliche Kompetenzen:

\begin{center}
\begin{tabulary}{\linewidth}{lclp{5cm}}
\toprule
\E{Gemeinde} 
 &
 ${\rightarrow}$
 &
Bürgermeister
&
\TakeNotesMedium
\\
\E{Bezirk} 
 &
 ${\rightarrow}$
 &
Bezirkshauptmann
&
\TakeNotesMedium
\\
\E{Bundesland}
 &
 ${\rightarrow}$
 &
Landeshauptmann
&
\TakeNotesMedium
\\
\E{Bundesgebiet}
 &
 ${\rightarrow}$
 &
Innenminister%, Bundeskanzler 
&
\TakeNotesMedium
\\\bottomrule
\end{tabulary}
\end{center}

\subsubsection{Die Katastrophenwarnung mittels Sirene alarmiert die
Bevölkerung}
\label{topic:alarmierung}

\includegraphics[width=\linewidth]{./images/noauth/khd-sirenen.png}

\subsection{Katastrophenschutz in Wien}
\label{topic:katastrophenschutz}
\TODO{EMHO}{2010-08-04}{Katastrophenhilfseinheiten fehlen!}{}
\begin{center}
\begin{tabular}{p{4.3190002cm}m{8.981cm}}
Wien ist für den Groß\-schadens- und Katastrophenfall gut
gerüstet. Tausende Helfer sind dazu im Katastrophen\-schutz-Kreis
(\emph{{\quotedblbase}K-Kreis{\textquotedblleft}}) org\-anisiert.

Der K-Kreis ist ein Zu\-sammen\-schluss der beruflichen und freiwilligen
Hilfs- und Ein\-satz\-organisationen sowie wichtiger Behörden.
Dreh\-scheibe und gemeinsame Anlaufstelle ist dabei die Stadt Wien,
Magistrats\-direktion - Katastrophen\-manage\-ment und
Sofort\-ma{\ss}\-nahmen und die Helfer Wiens.
 &
\includegraphics[width=9.046cm,height=13.623cm]{./images/noauth/k-kreis08.png}
\CaptionofDiff{figure}{Der K-Kreis.
\SourcePicture{http://www.diehelferwiens.at/k-kreis , Copyright: Die Helfer
Wiens, Grafik: Erdle/Zimmermann/Schimanek}}{Der K-Kreis}
\\
\end{tabular}
\end{center}

\subsubsection{Katastrophenhilfe}

Hilfsmaßnahmen die von uns als Sanitätsorganisation erwartet
werden:

\begin{itemize}
\item Bergen und Retten, Versorgung und Transport 
\item Pflegerische Betreuung 
\item Sorge für Unterkunft und Verpflegung 
\item Karitative Betreuung von Hilfsbedürftigen
\end{itemize}

\begin{center}
\begin{tabularx}{\linewidth}{lX}
\toprule
\TabHeading{Sanitätseinsatz}
 &
\begin{itemize}
     \item Retten/Bergen
     \item Triage
     \item Lebensrettende Sofortmaßnahmen
     \item Erstversorgung
     \item Sanitätshilfe
     \item Abtransport
\end{itemize}
\\\midrule
\TabHeading{Pflegeeinsatz}
&
 Wenn die pflegerische Versorgung durch die bestehenden Einrichtungen
 nicht mehr sichergestellt ist.
 \begin{itemize}
     \item Erweiterung der vorhandenen Behandlungsmöglichkeiten
     \item Errichtung von Hilfskrankenhäusern
 \end{itemize}
\\\midrule
\TabHeading{Sozialeinsatz}
&
Wenn Menschen infolge einer Katastrophe ihre
Selbstversorgungsmöglichkeit verloren haben.

\begin{itemize}
    \item Einrichtung/Betrieb von Notunterkünften/Lagern
    \item Beschaffung, Zubereitung und Ausgabe von Lebensmitteln, Kleidern,
    Gebrauchsgegenständen
    \item Fürsorgliche Hilfe, Registrierung und Betreuung
\end{itemize}
\\\bottomrule
\end{tabularx}
\end{center}


\subsection{Grundsätze des Einsatzes}
\label{topic:einsatzgrundsaetze}

\Fazit{Grundkomponenten: Führung -- Zeit -- Raum} 

Ziel ist es ohne überstürztes Handeln die Ausdehnung der Krisenlage
auf benachbarte und andere Räume zu verhindern. Dabei wird
gemäß der gesetzlichen Bestimmungen von Kräften der
Einsatzorganisationen, Behörden und anderen Einrichtungen koordiniert
vorgegangen.

\subsection{Führung}
\label{topic:katastrophe-fuehrung}
\Definition{Steuerndes Einwirken auf andere Menschen um ein bestimmtes Ziel zu
erreichen. }

Die Führung ist ein unentbehrliches Steuerungsinstrument zur
Erfüllung des Einsatzauftrages
({\quotedblbase}Ziel{\textquotedblleft}) und dem richtigen Einsatz von
Personal, Material und Transportmitteln 

Im Schadensraum gilt: {\quotedblbase}\Es{Ein einziger
Verantwortlicher}{\textquotedblleft} mit Entscheidungsbefugnis und den
nötigen Führungsmittel \E{({\quotedblbase}Einheit der
Führung{\textquotedblleft})}. Dieser wirkt mittels möglichst
einfacher Ma{\ss}nahmen \E{({\quotedblbase}Einfachheit{\textquotedblleft})}
auf den Einsatzverlauf ein
\E{({\quotedblbase}Handlungsfreiheit{\textquotedblleft})}. Mittel und
Mannschaften werden dabei gemäß ihrer Eigenart und
Leistungsfähigkeit eingesetzt \E{({\quotedblbase}\"Okonomie der
Kräfte{\textquotedblleft})}. Er muss rasch auf \"Anderungen der
Situation reagieren \E{({\quotedblbase}Beweglichkeit{\textquotedblleft})}
und gleichzeitig vorausschauend handeln
\E{({\quotedblbase}Reservenbildung{\textquotedblleft})}.

\Cave{{\quotedblbase}Führen{\textquotedblleft} und
{\quotedblbase}Behandeln{\textquotedblleft} ist nicht gleichzeitig
möglich!}

\paragraph{Sanitäter und Notärzte müssen Führungsaufgaben
übernehmen}~

Der erste am Schadensort eintreffende Sanitäter muss:

\begin{enumerate}
    \item Lagemeldung abgeben
    \item mit bereits anwesenden Einsatzkräften die ersten Maßnahmen
    zur Rettung und Behandlung festlegen
    \item festlegen, wohin die Patienten im Schadensraum gebracht werden 
    sollen (Beginn einer Triage)
    \item festlegen, der Grundstrukturen einer SanHiSt (Wagenhalteplatz, etc.)
\end{enumerate}

\subsection{Zeit}
\label{topic:katastrophe-phasen}

Jedes Großschadens- oder Katastrophenereignis zeigt unabhängig von seiner
Ursache i.d.R. folgenden Zeitablauf:

\begin{table}[H]
\Captionof{table}{Phasen des Einsatzes.}

\begin{tabularx}{\linewidth}{XXX}
\TabHeading{1. Phase}
 &
\TabHeading{2. Phase}
 &
\TabHeading{3. Phase} 
\\\midrule\addlinespace
Isolation (Minuten)
 &
Retten (Stunden)
 &
Wiederherstellung (Tage)
\\\addlinespace
Schaden begrenzen  

- Lage erkennen 

- Melden
 &
Ordnung schaffen 

- Kommandoposten
 &
Lokale Einsatzleitung 

- Verbindung 

- Lage beurteilen
\\\addlinespace
\"Uberleben
 &
Weiterleben
 &
Endbehandlung
\\\addlinespace
Laien: 

- Nothilfe/Selbsthilfe 

- Retten / Erste Hilfe
 &
Geschulte Helfer: 

- LRSM 

- Transportfähigkeit
 &
Spezialisten: 

- Spital nach Dringlichkeit behandeln
\\\addlinespace
\multicolumn{3}{c}{\includegraphics[width=0.5\linewidth,height=1em]{./images/khd/arrow-right.png}}
\\\addlinespace\bottomrule
\end{tabularx}
\end{table}


\section[Sanitätsdienstliche Organisation im
Schadensraum und in der SanHiSt]{Hilfe braucht Raum: Sanitätsdienstliche Organisation im
Schadensraum und in der Sanitätshilfsstelle (SanHiSt)}
\label{topic:sanhist}
% TODO : Layout-Anpassung
\TODO{GABS}{yyyy-mm-dd}{Layout-Anpassung notwendig.}{Dieser Teil entspricht nicht dem Standard-Layout!}

\Layout{0}{Übersicht}{
Der Einsatz allgemein wird unterteilt in: 

\begin{itemize}
  \item Schadensraum
  \item Transportraum
  \item Hospitalisationsraum
\end{itemize}

\begin{description}
  \item[Der Hospitalisationsraum] umfasst die Transportziele und
  Aufnahmespitäler.
  \item[Der Transportraum] umfasst in diesem Zusammenhang alle Mittel und Wege um
  die Patienten vom Schadens\-raum zum Hospitalisationsraum zu
  verbringen. Dieser {\quotedblbase}allgemeine
  Transportraum{\textquotedblleft} ist nicht zu ver\-wechseln mit dem
  Transportraum der Sanitätshilfsstelle (s.u.).
  \item[Der Schadensraum] ist der Ort des eigentlichen (Schadens-)Geschehens. Es
  gibt in ihm eine einheitliche Organisations\-form für den
  Sanitätsdienst: Die \T{Sanitätshilfsstelle} (\T{SanHiSt}), welche aus den
  drei Räumen \T{Triageraum}, \T{Behandlungsraum} und
   \T{Transportraum} besteht. Definiert wird Sie z.B. in der ASB\"O Rahmen\-vor\-schrift für
  Großunfälle und den Durchführungsbestimmungen für
  Großunfälle.
\end{description}
}{
\begin{itemize}
  \item Schadensraum $\rightarrow$ SanHiSt
  \begin{itemize}
    \item Triageraum
    \item Behandlungsraum
    \item Transportraum
  \end{itemize}
  \item Transportraum
  \item Hospitalisationsraum
\end{itemize}
}{}{}{}


 



\begin{minipage}{12.605cm}
\includegraphics[width=\linewidth]{./images/khd/sanhist-plan-edited.png}
\Captionof{figure}{Gliederung des Einsatzraumes}
\begin{multicols}{2}

\begin{description}
    \item[\T{\"Außere Absperrung}:] Großräumige Verkehrsumleitung
    \item[\T{Innere Absperrung}:] hält Schaulustige fern, hält Personen in Panik
    zurück
    \item[\T{Sicherheitsring}:] Absperrung des Schadenplatzes, Sicherheit für
    Einsatzkräfte
    \item[\T{Pforte}:] Einbahnsystem für Einsatz-Kfz, keine Fremdfahrzeuge,
    möglichst nur eine Pforte
    \item[\T{Leiter Information}] (i) für Weiterleitung von Informationen an die
    Außenwelt (Presse, Angehörige, ...)
    \item[\T{Bergetriage}] (BT)
    \item[\T{Sanitätshilfsstelle} (\T{SanHiSt}):] mit Triageraum, Behandlungsraum und
    Transportraum;
    \item[\T{Einsatzleiter}] (EL)
    \item[\T{Leitender Notarzt}] (LNA)
    \item[\T{Mobile Leitstelle}] (MLS)
    \item[\T{Material- und Meldestelle}] (M)
    \item[\T{Sammelstelle Tote}] (t) und
    \item[\T{Sammelstelle Unverletzte}] (U)
\end{description}


\end{multicols}
\end{minipage}

\clearpage
\subsection{Triageraum}

\Layout{60}{Bergetriage}{%
Die Bergetriage nimmt die Festlegung der Bergepriorität vor, d.\,h. in
welcher Reihenfolge die Patienten zur Triagestelle gebracht werden sollen.
Sie wird nur gebildet wenn keine Gefahrenzone und genügend Personal vorhanden
ist.

Das Team muss mind. aus einem NFS und einem NA bestehen. 
Es erfolgt eine Sichtung der Betroffenen und Festlegung 
der Behandlungsdringlichkeit. Die Bergetriage behandelt nicht!
}{
\begin{itemize}
    \item Festlegung der Bergepriorität.
    \item Mind. NA + NFS.
    \item Gefahren beachten.
    \item Sichtung und Festlegung der
    Behandlungsdringlichkeit.
    \item Nicht behandeln!
\end{itemize}
}
{./images/noauth/sanhist-bergetriage.png}
{}{}





%%%%%%%%%%%%%%%%%%%%%%%%%%%%%%%%%%%%%%%%%%%%%%%%%%%%%%%%%%%
% TODO Dynamik I-IV: Überarbeitung notwendig.
\TODO{GABS}{yyyy-mm-dd}{Dynamik I-IV: Überarbeitung notwendig.}{Dynamik I-IV: Überarbeitung notwendig. Aussage der i-IV-Dynamik kommt nicht
   gut herus.}
%%%%%%%%%%%%%%%%%%%%%%%%%%%%%%%%%%%%%%%%%%%%%%%%%%%%%%%%%%%

\Layout{0}{Triage (Sichtung)}{%
Aufgabe der Triagestelle ist die Zuordnung der Patienten zu den Triagegruppen
(\begin{inparadesc}
\item[\Code{I~~}] Sofortbehandlung, lebensrettende Noteingriffe; 
\item[\Code{II~}] Dringende Behandlung, Schwerverletzte; 
\item[\Code{III}] Minimalbehandlung; 
\item[\Code{IV~}] Abwartende Behandlung
\end{inparadesc}) und den Sammelstellen. \E{Spätestens} hier erfolgt die
zentrale Registrierung aller Betroffenen mittels Patientenleitsystem (PLS). 

Die Triage erfolgt schnell, mit geringem Aufwand und unter hohem 
Zeitdruck. Der maximale Zeitaufwand pro Patient beträgt bei gehenden Patienten
1 Minute und bei liegenden Patienten 3 Minuten. Es erfolgt \E{keine
detaillierte} Diagnostik.

\begin{center}
% \includegraphics[width=2.cm]{./images/noauth/sanhist-triage.png}
\TODO{GABS}{2010-04-27}{Foto Triageplatz}{}
\end{center}

Es werden zwei Triageplätze eingerichtet, damit sich der Notarzt nur zum
nächsten Patienten umdrehen braucht und abwechselnd triagieren kann. 
Die Vorbereitung des Patienten erfolgt durch Sanitätsfachpersonal und umfasst
das Entkleiden sowie die Vornahme des Ersteinschätzungsblockes. 

}{
\begin{itemize}
        \item Zentrale Registrierung aller Patienten mittels Patientenleitsystem
    (PLS)
        \item Zuordnung der Patienten zu Triagegruppen I\,--\,IV und
        Sammelstellen U und t
    \item Schnell, mit geringem Aufwand, hoher Zeitdruck.

  Maximaler Zeitaufwand pro Patient gesamt:
    \begin{itemize}
        \item Gehender 1 Minute
        \item Liegender 3 Minuten
    \end{itemize}    
    \item Keine detaillierte Diagnostik 
\end{itemize}
}
{./images/noauth/sanhist-tragestelle-pic.jpg}
{}
{}



\Layout{0}{Exkurs: Triagegruppen}{%
Die Zuordnung zu einer Triagegruppe kann sich im Verlauf verändern,
eine ständige Nachtriage ist erforderlich! Änderungen ergeben sich besonders durch:

\begin{enumerate}
    \item Verschlechterung des Patientenzustandes: z.\,B. III → I.
    \item Verbesserung des Patientenzustandes
    (z.\,B. nach erfolgter Behandlung): z.\,B. I → III.
    \item Verbesserung der Versorgungslage (zusätzliches 
    Personal bzw. Material verfügbar, oder Behandlungskapazitäten sind frei
    geworden): z.\,B. IV → I.
\end{enumerate}
}{

}{}{}{}








\subsection{Behandlungsraum}



\Layout{60}{Behandlungsstellen}{%
Es werden vier Behandlungsstellen entsprechend der Triagegruppen eingerichtet.
\prettyref{tab:khd-behandlungsstellen-uebersicht} gibt eine Übersicht über
deren Aufgaben und Eigenschaften. 

\vspace{2cm}

}{

}
{./images/noauth/sanhist-behandlung.png}
{Symbol Behandlungsraum.}
{\Abk{BMI}}




% \begin{tabulary}{\linewidth}{lL}
% \toprule
% \begin{minipage}{3.5cm}
% \includegraphics[width=1.489cm,height=1.61cm]{./images/noauth/sanhist-behandlung.png}
% \end{minipage}
%  &
% Es werden drei Behandlungsstellen eingerichtet
% \\
% \begin{minipage}{3.5cm}
% \includegraphics[width=1.629cm,height=1.489cm]{./images/noauth/sanhist-triage-1}
% \end{minipage}
% 
% 
%  &
% Sofortbehandlung vor Ort
% 
% T1 Lebensrettende Soforteingriffe 
% \\
% \begin{minipage}{3.5cm}
% \includegraphics[width=1.669cm,height=1.578cm]{./images/noauth/sanhist-triage-2.png}
% \end{minipage}
%  &
% Dringende Behandlung
% 
% T2 Schwerverletzte
% \\
% \begin{minipage}{3.5cm}
% \includegraphics[width=1.601cm,height=1.587cm]{./images/noauth/sanhist-triage-3.png}
% \includegraphics[width=1.587cm,height=1.641cm]{./images/noauth/sanhist-triage-4.png}
% \end{minipage}
% 
%  &
% Warten
% 
% T3 Warten Leichtverletzte
% 
% T4 Warten Schwerverletzte
% \\\bottomrule
% \end{tabulary}



















% 
% 
% \subsubsection{I -- Sofortbehandlung vor Ort}
% 
% Alle Patienten, welche ohne sofortige Behandlung sterben bzw. schwere
% Schäden erleiden
% 
% \begin{flushleft}
% \begin{tabular}{m{5.642cm}m{4.379cm}m{3.076cm}}
% Welche Patienten?
% 
% \begin{itemize}
%     \item Atemstörungen
%     \item Spannungspneumothorax
%     \item schwere äu{\ss}ere Blutungen
%     \item schwerer traumatischer Schock
% \end{itemize}
% &
% Ma{\ss}nahmen
% 
% \begin{itemize}
%     \item \"Uberprüfung der Triage
%     \item Volumenersatz
%     \item Intubation
%     \item Thoraxdrainage
%     \item Noteingriffe
% \end{itemize}
% &
% \includegraphics[width=3cm]{./images/noauth/sanhist-triage-1.png}
% \\
% \end{tabular}
% \end{flushleft}
% 
% 
% \subsubsection{II -- Dringende Behandlung}
% 
% Behandlung und Herstellung der Transportfähigkeit von schwer
% verletzten/erkrankten Personen, eine weitere Behandlung im Krankenhaus
% ist i.d.R. notwendig.
% 
% \begin{flushleft}
% \begin{tabular}{m{5.642cm}m{4.379cm}m{3.076cm}}
% Welche Patienten?
% 
% keine unmittelbare Lebensbedrohung
% 
% \begin{itemize}
%     \item SHT
%     \item Wirbelverletzungen mit Lähmungen
%     \item Bauchtrauma, innere Blutungen
%     \item schwere Thoraxtraumen
%     \item gro{\ss}e Gefä{\ss}verletzungen
%     \item offene Knochenbrüche und  Gelenks\-verletzungen
%     \item Augenverletzungen
%     \item gro{\ss}e Weichteilverletzungen
%     \item Verbrennungen 15 - 30\%
%     \item geschlossene Knochenbrüche und Gelenks\-ver\-letzungen
% \end{itemize}
% &
% Ma{\ss}nahmen
% 
% \begin{itemize}
%     \item Verbände
%     \item Schockbehandlung
%     \item Fixierung/Schienung
%     \item Schmerzbekämpfung
%     \item pflegerische Ma{\ss}nahmen
%     \item LRSM wenn nötig
% \end{itemize}
% &
% \includegraphics[width=3cm]{./images/noauth/sanhist-triage-2.png}
% 
% \\
% \end{tabular}
% \end{flushleft}
% 
% 
% \subsubsection{III -- Warten Leichtverletzte}
% 
% Alle Patienten, welche nur einer Minimalbehandlung bedürfen.
% 
% \begin{flushleft}
% \begin{tabular}{m{5.642cm}m{4.379cm}m{3.076cm}}
% Welche Patienten?
% 
% \begin{itemize}
%     \item Verbrennungen bis 15\%
%     \item kleine Weichteilwunden
%     \item einfache Knochenbrüche
%     \item Prellungen, Zerrungen
%     \item leichtes SHT
% \end{itemize}
% &
% Ma{\ss}nahmen
% 
% \begin{itemize}
%     \item \"Uberprüfung der Triage
%     \item pflegerische Ma{\ss}nahmen
%     \item Betreuung (Kriseninterventionsteam (\Abk{KIT})
%     \item Entlassung ambulant Behandelter zur SaStU
% \end{itemize}
% &
% \includegraphics[width=3cm]{./images/noauth/sanhist-triage-3.png}
% \\
% \end{tabular}
% \end{flushleft}
% 
% 
% \subsubsection{IV -- Warten Schwerverletzte}
% 
% Alle Patienten, welche mit dem zur Verfügung stehenden Mitteln
% (Personal, Material) noch nicht versorgt werden können . Abwartende
% Behandlung.
% 
% \begin{center}
% \begin{tabular}{m{5.702cm}m{4.0550003cm}m{3.3439999cm}}
% Welche Patienten?
% 
% \begin{itemize}
%     \item schweres Polytrauma
%     \item schwerstes \Abk{SHT}
%     \item offene Körperhöhlentraumen
%     \item Mehrhöhlenverletzungen
%     \item Verbrennungen über 50\%
% \end{itemize}
% &
% Ma{\ss}nahmen
% 
% \begin{itemize}
%     \item \"Uberprüfung der Triage
%     \item Behandlung  (Schmerzbekämpfung)
%     \item Psychologische Betreuung, evtl. durch \Abk{KIT}
%     \item \"Uberwachung
% \end{itemize}
%  &
% \includegraphics[width=3cm]{./images/noauth/sanhist-triage-4.png}
% \\
% \end{tabular}
% \end{center}









\subsection{Transportraum}

\Layout{60}{Verladestelle}{%
Die Verladestelle sollte -- sofern möglich -- in der Nähe der Behandlungsstelle
II errichtet werden. Es muss protokolliert werden, welcher Patient mit welchem
Transportmittel in welches Krankenhaus eingeliefert wird.

Es gibt eine definierte Ein- und Ausfahrt, eine Einbahnregelung ist tunlichst
anzustreben! }{

}
{./images/noauth/sanhist-verladestelle.png}
{}
{}





\Layout{60}{Wagenhalte\-platz}{%
Der Wagenhalteplatz muss \Es{frühzeitig} bestimmt werden, da weitere Kräfte
rasch nachfolgen werden. Er sollte in ausreichender \Es{Entfernung} festgelegt
werden um den späteren Aufbau der weiteren SanHiSt nicht zu behindern. Bedenke
dabei, dass der Triage- und Behandlungsraum (und die Verladestelle) 
viel Platz brauchen! Der Wagenhalteplatz selber muss allerdings auch groß genug
angelegt sein. Auch die Witterung und die Beschaffenheit des Untergrunds ist zu
beachten.

Wenn
möglich herrscht Schrägparkordnung,  um ein freies Wegfahren zu gewährleisten.
Die Fahrzeuge sollen nach Einsatzmittelkategorien (\Abk{NAW}, \Abk{RTW}, \ldots)sortiert abgestellt werden. Die Schlüssel müssen stecken gelassen werden!
}{
\begin{itemize}
    \item Frühzeitig, Abstand von SanHiSt. 
    \item Groß genug; Bodenbeschaffenheit und Witterung bedenken!
    \item Schrägparkordnung, freies Heck zum Ent-/Beladen. 
    \item Sortiert nach Einsatzmittel. 
    \item Schlüssel stecken lassen. 
\end{itemize}
}
{./images/noauth/sanhist-wagenhalteplatz.png}
{}
{}









%  Falscher Wagenhalteplatz
% 
% 
% Warning: Image ignored
% Unhandled or unsupported graphics:
%\includegraphics[width=3.375cm,height=2.438cm]{AASdev20090228-img129}




\Layout{0}{Hubschrauberlandestelle}{%
\begin{inparaitem}
    \item Die allgemeinen Vorschriften für Hubschrauberlandeplätze sind zu beachten
    \item Freie ebene Fläche
    \item Sichere Umgebung (Stromleitungen!)
    \item Weit genug weg vom Behandlungsraum (Luftwirbel!)
\end{inparaitem}
% TODO BILD: Hubschrauberlandestelle
}{

}{}{}{}







\Layout{22}{Exkurs: Transportmanagement}{%
\begin{inparaitem}
    \item Spiralförmige Verteilung
    \item Zeitliche Staffelung
    \item Fachliche Zuordnung
    \item Katastrophe nicht ins Spital verlagern!
\end{inparaitem}
}{

}
{./images/noauth/khd-transportmanagement.png}
{}
{}






\subsection{Sammelstellen}

\Layout{60}{Sammelstelle Unverletzte}{%
Ziel: Betreuung Unverletzter

\begin{inparaitem}
    \item Betreuen
    \item Beruhigen
    \item Informieren
    \item \Abk{KIT}\footnote{In Wien ist die Akutbetreuung Wien (\Abk{ABW})
    zuständig.} integrieren
\end{inparaitem}

\vspace{1cm}
}{

}
{./images/noauth/sanhist-sammelstelle-unverletzte.png}
{}
{}




\Layout{60}{Sammelstelle Tote}{%
Sollte diskret eingerichtet werden.

Abgeschirmt von der Sammel\-stelle Unverletzte, dem Be\-handlungs\-raum
und der Außen\-welt

\vspace{1.5cm}
}{

}
{./images/noauth/sanhist-sammelstelle-tote.png}
{}
{}







\subsection{Organisation}
\label{topic:khd-organisation}


\Layout{60}{Material- / Meldestelle}{%
Sie dient zur Registrierung des Personals und Materials.  Eintreffende
Einheiten haben sich bei dieser zu melden.
Sie steht im engen Kontakt zum Leiter SanHiSt und EL. 

\vspace{1.5cm}

}{

}
{./images/noauth/sanhist-material-meldestelle.png}
{}
{}




\Layout{60}{Mobile Leitstelle (MLS)}{%
Die \Abk{MLS} beherbergt i.d.R. die Ein\-satz\-leitung sowie wichtige
    Administrations-, Tele\-kom\-munikations- und Funkeinrichtungen.
   Bei Bedarf werden u.U. hier Handfunkgeräte an die Bereichsleiter
    ausgegeben.
    Die MLS ist die Brücke zwischen verschiedenen Funkkanälen und
    zuständig.
     Wenn die MLS die Einsatzleitung beinhaltet wird sie mit einem
    rotem Dreh- oder Blinklicht markiert.

}{

}
{./images/khd/khd-mls.jpg}
{}
{}




\subsection{Führungsstruktur in der SanHiSt}

\begin{center}
\begin{tabular}{m{1.8cm}m{4.5cm}|m{1.8cm}m{4.5cm}}
\includegraphics[width=1.654cm,height=1.624cm]{./images/noauth/sanhist-einsatzleiter.png}
&
\Es{Einsatzleiter} (Rettungsdienst) 

Er ist der Vorgesetzte für die
eingesetzten Sanitäter und sonstigen sanitätsdienstlichen Hilfskräfte. Er hat
\E{keine} direkte Weisungsbefugnis gegenüber anderen eingesetzten Organisationen. Er kann Ärzten
\E{keine Weisungen in medizinischen Fragen} geben. 
&
% Unhandled or unsupported graphics:
\includegraphics[width=1.62cm,height=1.62cm]{./images/noauth/sanhist-material-meldestelle.png}
&
\Es{Leiter Material-/Meldestelle}
\\ 
% Unhandled or unsupported graphics:
\includegraphics[width=1.62cm,height=1.62cm]{./images/noauth/sanhist-leiter-sanhist.png}
&
\Es{Leiter Sanhist}
&
% Unhandled or unsupported graphics:
\includegraphics[width=1.489cm,height=1.61cm]{./images/noauth/sanhist-behandlung.png}
&
\Es{Leiter Behandlung}
\\ % Unhandled or unsupported graphics:
\includegraphics[width=1.62cm,height=1.62cm]{./images/noauth/sanhist-lna.png}
&
\Es{LNA Leitender Notarzt} 

Der LNA ist der Vorgesetzte der eingesetzten Ärzte.
& 
% Unhandled or unsupported graphics:
\includegraphics[width=1.62cm,height=1.62cm]{./images/noauth/sanhist-leiter-transport.png}
&
\Es{Leiter Transport}
\\ 
% Unhandled or unsupported graphics:
\includegraphics[width=1.62cm,height=1.62cm]{./images/noauth/sanhist-leiter-triage.png}
&
\Es{Leiter Triageraum/Triage}
& 
% Unhandled or unsupported graphics:
\includegraphics[width=1.62cm,height=1.62cm]{./images/noauth/sanhist-information.png}
&
\Es{Leiter Information}
\\
\end{tabular}
\end{center}


\section{Patientenleitsystem (PLS)}
\label{topic:pls}
% TODO : Layout-Anpassung
\TODO{GABS}{yyyy-mm-dd}{Layout-Anpassung notwendig.}{Dieser Teil entspricht nicht dem Standard-Layout!}

\Layout{0}{\TxBeschreibung}{
Das \index{Patientenleitsystem}\T{Patientenleitsystem} (\index{PLS}\T{PLS})
ermöglicht eine \Es{eindeutige Kennzeichnung} von Patienten und Zuordnung
von Behandlungsmaßnahmen und Befunden. Der wichtigste Bestandteil
ist die \index{Patientenleittasche}\T{Patientenleittasche}
(\index{PLT}\T{PLT}). 

Sie besteht aus einer orangenen Hülle welche \E{eindeutige nummeriert}
ist. Auf dieser Kunststoffhülle können wichtige Notizen wie
Triagegruppe, etc., vorgenommen werden. In der Hülle befinden sich
das \E{Behandlungsprotokoll} (blau), \E{Identifikationsprotokoll} (rosa),
\E{Zusatzkarten}, \E{Kontaminationsaufkleber} und ein Bogen mit gleichlautend
\E{nummerierten Klebeetiketten}. Am unteren Ende befinden sich \E{zwei
Abschnitte} welche die Rückverfolgbarkeit des Transports
sicherstellen. Die PLT ist mit einer Gummischnur versehen damit sie dem
Patienten umgehängt werden kann. 

\Es{\E{Spätestens} an der Triagestelle} bekommt jeder Betroffene {\quotedblbase}seine{\textquotedblleft} \Abk{PLT}. Auch Unverletzte
und Verstorbene werden mittels \Abk{PLS} registriert.

\paragraph{Zusatzkarten}

Bei einer eventuellen Bergetriage bekommt jeder Patient eine
Patientenleittasche (\Abk{PLT}). Patienten, welche \Es{bevorrangt gerettet}
werden sollen, werden zusätzlich mit der gelben
{\fcolorbox{Black}{Yellow}{\sffamily\bfseries DRINGEND}}\,-Karte
gekennzeichnet. Die Karte wird auf der Triagestelle wieder entfernt. Tote
erhalten zusätzlich zur \Abk{PLT} eine schwarze
{\fcolorbox{Black}{White}{\sffamily\bfseries VERSTORBEN}}\,-Karte.

\paragraph{Behandlungsprotokoll / Identifikationsprotokoll}

Die Protokolle werden im Behandlungsraum ausgefüllt wenn Zeit dafür
ist. Das Ausfüllen der Protokolle darf den Abtransport unter keinen
Umständen verzögern. Die Protokolle können auch leer bleiben.

\paragraph{Kontaminationsaufkleber}\index{Kontaminationsaufkleber}

Wenn der Patient kontaminiert ist, wird ein Kontaminationsaufkleber in
Form eines gelben Dreiecks mit schwarzem Rand auf die \Abk{PLT} geklebt. Dies soll
alle nachfolgenden Personen, welche mit dem Patienten in
Kontakt treten, warnen. Die Aufkleber sind auf \Es{beiden
Seiten} der \Abk{PLT} anzubringen!

\paragraph{PLT-Nummer}

Die \Abk{PLT} und damit der Patient ist durch eine \Es{eindeutige, einmalige
Nummer} \Es{identifiziert}. Mittels \Es{Selbstklebeetiketten} können die
Nummern der Patienten in verschiedenen Listen schnell erfasst werden. Das ist
in vielen Bereichen besonders wichtig (Triage, Behandlung, Transport). Außerdem
können alle Effekte des Patienten damit gekennzeichnet werden.

Die Patientenleittasche verbleibt solange am Patienten, \E{bis das
Aufnahmeverfahren im Krankenhaus abgeschlossen ist}, d.h. bis der Patient im
Spital administrativ erfasst ist. Anschließend wird die \Abk{PLT} der
Krankengeschichte des Patienten beigefügt.

\Fazit{Die Nummerierung der \Abk{PLT} ermöglicht die Zuordnung und Erfassung von
Patienten. Das ist eine der Hauptaufgaben der \Abk{PLT}!}

\A{PLT-Fotos \ldots}


\includegraphics[width=\linewidth]{./images/khd/plt-inhalt-beschriftet.jpg}

\begin{multicols}{2}

\begin{description}
    \item[a] Zusatzkarte {\fcolorbox{Black}{White}{\sffamily\bfseries
VERSTORBEN}}
    \item[b] Zusatzkarte 
{\fcolorbox{Black}{Yellow}{\sffamily\bfseries DRINGEND}}
    \item[c] Klebeetiketten mit PLT-Nummer
    \item[d] Aufkleber für kontaminierte
Gegenstände 
    \item[e] Behandlungsprotokoll (blau)
    \item[f] Identifikationsprotokoll (rosa)
\end{description}

\end{multicols}
}{
\begin{itemize}
  \item Patientenleitsystem (PLS)
  \item Patientenleittasche (PLT)
  \begin{itemize}
    \item Behandlungsprotokoll (blau)
    \item Identifikationsprotokoll (rosa)
    \item Zusatzkarten
    \begin{itemize}
      \item Dringend
      \item Verstorben
    \end{itemize}
    \item Kontaminationsaufkleber
    \item Klebeetiketten mit PLT-Nummer
  \end{itemize}
\end{itemize}
}{}{}{}
 
\subsubsection{PLT Vorderseite}

\noindent\Parallel{
\paragraph{Feld \emph{"`Triage"'}} \textcircled{a}

Kennzeichnung der Triagegruppe:

Bei Triage nur Einteilung in {\sffamily I}, {\sffamily II},
{\sffamily III} oder {\sffamily IV}

\paragraph{Zweiter Triageblock} \textcircled{b}

Dient zur 
\begin{itemize}
    \item Umtriage oder
    \item Spitalstriage
\end{itemize}

\paragraph{Feld \emph{"`Transport"'}} \textcircled{c}

Kennzeichnung der Transportpriorität, wird später
entschieden

\begin{description}
    \item[{\sffamily A}: Hohe Transportpriorität:] Der
    rasche Transport zur frühzeitigen Fachbehandlung nach
    ärztlicher Notversorgung ist angesagt.
    \item[{\sffamily B}: Niedrige Transportpriorität:] Der Transport
zur Fachbehandlung kann verzögert erfolgen. 
\end{description}

  

 




}{
\includegraphics[width=\linewidth]{./images/khd/plt-vorderseite-beschriftet-50pc.jpg}
}

%%%%%%%%%%%%


\begin{minipage}[t]{0.48\linewidth}
\includegraphics[width=\linewidth]{./images/khd/plt-vorderseite-umtriagierung15-50pc.jpg}
\Captionof{figure}{PLT: Umtriagierung innerhalb der ersten
Viertelstunde}
\begin{itemize}
\item Mittels Pfeil von einem zum nun zutreffenden Triagefeld
\item In der gleichen Zeile
\end{itemize}
\end{minipage}
\hspace{0.04\linewidth}
\begin{minipage}[t]{0.48\linewidth}
\includegraphics[width=\linewidth]{./images/khd/plt-vorderseite-umtriagierung60-50pc.jpg}
\Captionof{figure}{PLT: Umtriagierung nach längerer Zeit}
\begin{itemize}
    \item Neues Kreuz beim aktuellen Triagefeld
    \item in der unteren Zeile
\end{itemize}
\end{minipage}
	





\noindent\Parallel{
\paragraph{Rückseite}~

Der Triagearzt trägt auf der Rückseite die zu setzenden
\Es{Ma{\ss}nahmen} ein:

\begin{itemize}
    \item Atmung ($O_2$, Intubation, usw..)
    \item Kreislauf (Blutstillung usw...)
    \item Medikamente
    \item Lagerung
\end{itemize}

Der Durchführende (Arzt, NFS, RS) bestätigt die Durchführung der
Ma{\ss}nahmen im rechten Feld durch eine Paraphe.

\vspace{10em}

\paragraph{Abriss für das Zielspital:}~

Wird bei der Einlieferung mit der Ankunftszeit versehen und verbleibt in
der Krankengeschichte.

\paragraph{Abriss für die SanHiSt -- Leiter Transport:}

Eintragen: Zielspital, Kfz Nummer, Verladezeit

Verbleibt beim Leiter Transport/Protokollführer
}{
\includegraphics[width=\linewidth]{./images/khd/plt-rueckseite-angekreuzt-50pc.jpg}
}








